\begin{proof}[Proof of \cref{obj:lem:cty-a1-a2-rectangle-angular-hypercube-all-orders}]
Throughout, a point of \(\mathbb R^2\) is written \(x=(x^{(1)},x^{(2)})\), and we abbreviate \(q=p+1\).  The arm sign attached to the side intervals \(I_1=[0,\infty)\), \(I_0=(-\infty,0)\) of \cref{obj:def:cty-a1-a2-class} is written \((2t-1)\), as in the arm-mass clause there.

\emph{Two fixed smooth profiles.}
Fix once and for all a radially symmetric function \(\varphi:\mathbb R^2\to[0,1]\) of class \(C^\infty\) with
\[
\varphi(0)=1,
\qquad
\varphi(z)=0\ \text{ whenever }\ \lVert z\rVert_2\geq 1,
\]
let \(\varphi_{\mathrm r}:[0,\infty)\to[0,1]\) be its radial profile, so that \(\varphi(z)=\varphi_{\mathrm r}(\lVert z\rVert_2)\), and put
\[
b_k=\sup_{z\in\mathbb R^2}\lVert D^k\varphi(z)\rVert<\infty,
\qquad k\geq 0 .
\]
% lean: packingBump packingBump_zero packingBump_mem_Icc packingBump_eq_zero_of_one_le_norm packingBump_contDiff packingBumpDerivativeBound packingBump_iteratedFDeriv_le_derivativeBound
Fix also a nondecreasing \(C^\infty\) function \(\chi:\mathbb R\to[0,1]\) with \(\chi(s)=0\) for \(s\leq0\) and \(\chi(s)=1\) for \(s\geq1\).
% lean: angularCutoff angularCutoff_mem_Icc angularCutoff_eq_zero angularCutoff_eq_one

\emph{The constants.}
Set
\[
A=257+\bigl(1+q+1310720+\sum_{k=0}^{q}b_k\bigr),
\qquad
c=\frac1{24A},
\qquad
C=128,
\qquad
C_0=262144 .
\]
All four are positive, \(A>256\) and \(A\geq1\) because the bracket is positive, and, since every \(b_k\) is nonnegative and appears in the bracketed sum,
\[
b_k\leq A
\qquad\text{for every }0\leq k\leq q .
\]
% lean: causalHardConstructionBandwidthConstant packingBumpDerivativeScale packingBumpDerivativeBound_le_scale
Because \(2C=256<A\), the required order-zero separation \(2C<A\) holds; it in fact holds for every \(p\).
% lean: causalHardHypercubeConstructorData

\emph{The small-separation threshold.}
The map \(\Delta\mapsto A\Delta^{1/q}\) is continuous at \(0\) with value \(0\), so there is \(\delta_0>0\) such that every \(0<\Delta\leq\delta_0\) satisfies
\[
\Delta\leq\frac1{16},
\qquad
\Delta A\leq1,
\qquad
w:=A\Delta^{1/q}\leq\frac1{24}.
\]
% lean: causalHardPowerBandwidthEventuallySmall causalHardHypercubeConstructorData

\emph{The envelope threshold.}
The signed radial polynomial energy is coercive: there is \(c_p>0\) such that, for both \(t\in\{0,1\}\) and every \(v=(v_0,\ldots,v_p)^\top\in\mathbb R^{p+1}\),
\[
c_p\lVert v\rVert_2^2
\leq
\int_0^1\bigl\{v^\top r_p((2t-1)u)\bigr\}^2 u\,du ,
\]
with \(r_p\) the basis of \cref{obj:synth_71}.
% lean: signedRadialPolynomialEnergy_coercive
Put
\[
\gamma_p=\frac1{48}\cdot\frac\pi2\cdot c_p>0,
\qquad
L_0=\max\{48,\gamma_p^{-1}\},
\]
so that \(L_0\geq48\).
% lean: hardSquare_populationGramFloor_eventually causalHardA1A2Law_populationGram_certificate
Fix \(\nu\geq2\) and \(L\geq L_0\); in particular \(L\geq48\).  Fix also \(0<\Delta\leq\delta_0\).

\emph{Common geometry and grid.}
Take the fixed known geometry of \cref{obj:def:cty-known-geometry},
\[
\mathcal X=[-3,3]^2,
\quad
\mathcal A_1=[-1,1]\times[0,2],
\quad
\mathcal A_0=\mathcal X\setminus\mathcal A_1,
\quad
\mathcal B=\partial\mathcal A_1,
\quad
d(x,z)=\lVert x-z\rVert_2,
\]
with the uniform kernel \(K_\square\) of \cref{obj:def:cty-uniform-kernel}; the interface \(\mathcal B\) is the one displayed in \cref{obj:synth_35}.  With \(w=A\Delta^{1/q}\) put
\[
M=\left\lfloor\frac1{12w}\right\rfloor,
\qquad
x_j=\bigl(-\frac14+\frac{j}{2(M+1)},\ 0\bigr),
\quad j=1,\ldots,M,
\qquad
S_j=B_2(x_j,w),
\]
the balls being those of \cref{obj:synth_43}.
% lean: angularGridSize causalHardGridCenter causalHardCell
Each \(x_j\) has second coordinate \(0\) and first coordinate of modulus at most \(1/4\), so it lies on the middle of the bottom edge \(\{(s,0):|s|\leq1/2\}\) of the treated rectangle, hence on \(\mathcal B\).
% lean: causalHardGridCenter_mem_bottomEdge causalHardBottomEdge
Since \(w\leq1/24\) gives \(1/(12w)\geq2\) and therefore \(M\geq2\) and \(M\geq 1/(24w)\),
\[
M\geq\frac1{24w}=\frac1{24A}\Delta^{-1/q}=c\,\Delta^{-1/q}.
\]
% lean: angularGridSize_lower causalHardGrid_scaled_family_geometry
The centers are equispaced with gap \(1/\{2(M+1)\}\), so for \(i\neq j\)
\[
\lVert x_i-x_j\rVert_2=\frac{|i-j|}{2(M+1)}\geq\frac1{2(M+1)}\geq 4w\geq 2w,
\]
using \(2(M+1)\leq3M\leq 1/(4w)\), which is where \(M\geq2\) and \(M\leq1/(12w)\) are used.
% lean: causalHardGridCenter_dist causalHardGridCenter_separated
Consequently the disks \(S_j\) are pairwise disjoint, and each is contained in \(\mathcal X\) because \(|x_j^{(1)}|\leq1/4\) and \(w\leq1/24\).
% lean: causalHardGrid_scaled_family_geometry

\emph{The vertex laws.}
For \(\omega\in\{0,1\}^M\) let \(P_\omega\) be the law with the following two ingredients.  First, \(X\) has Lebesgue density \(f_{P_\omega}\) supported on \(\mathcal X\), given there by
\[
f_{P_\omega}(x)
=
\frac1{36}
\left\{
1+\sum_{j=1}^{M}\omega_j\,
t_\Delta\bigl(\lVert x-x_j\rVert_2\bigr)\,\eta_j(x)
\right\},
\]
where the radial tilt and the horizontal direction cosine at \(x_j\) are
\[
t_\Delta(r)
=
-\,\frac{2\Delta\,\varphi_{\mathrm r}(r/w)\,
\chi\bigl(\tfrac{r}{128\Delta}-1\bigr)}
{\max\{r/16,\ 8\Delta\}},
\qquad
\eta_j(x)=\frac{x^{(1)}-x_j^{(1)}}{\lVert x-x_j\rVert_2}
\ \ (x\neq x_j),
\quad
\eta_j(x_j)=0 .
\]
% lean: angularTilt angularRadialProfile packingDirectionCos packingAngularTerm packingAngularDensity causalHardScoreDensity causalHardScoreMeasure
Second, given \(X=x\) the two potential outcomes are conditionally independent Bernoulli variables whose means are the arm regressions of \cref{obj:def:cty-a1-a2-class},
\[
\mu_{0,P_\omega}(x)=\frac12,
\qquad
\mu_{1,P_\omega}(x)
=
\frac12+\frac{x^{(1)}}{16}
+\sum_{j=1}^{M}\omega_j\,\Delta\,
\varphi\!\left(\frac{x-x_j}{w}\right),
\]
both truncated to \([0,1]\) off \(\mathcal X\), with \(T=\mathbf 1\{X\in\mathcal A_1\}\) and \(Y=TY(1)+(1-T)Y(0)\).
% lean: causalHardControlProfile causalHardTreatmentProfile packingRegression packingAffineBaseline localizedPackingBump causalSelectedBernoulliKernel causalBernoulliPotentialOutcomeMeasure causalHardA1A2Law
Because the disks are \(4w\)-separated while \(\varphi\) vanishes outside the unit ball, at most one summand is nonzero at any \(x\); together with \(|x^{(1)}|/16\leq3/16\) and \(\Delta\leq1/16\) this gives
\[
\frac14\leq\mu_{1,P_\omega}(x)\leq\frac34,
\qquad x\in\mathcal X,
\]
so the truncation is inactive on \(\mathcal X\).
% lean: causalHardPackingRegression_mem_middleHalf causalHardTreatmentProfile_mem_middleHalf causalHardTreatmentProfile_eq_packingRegression_on_square
For the same reason the angular sum has modulus at most \(1/4\), whence
\[
\frac1{48}\leq f_{P_\omega}(x)\leq\frac5{144},
\qquad x\in\mathcal X.
\]
% lean: packingAngularDensity_mem_Icc causalHardScoreDensity_mem_Icc

\emph{Class membership.}
We check the clauses of \cref{obj:def:cty-a1-a2-class} in turn, using \(L\geq48\) throughout.

The support is the square \(\mathcal X=[-3,3]^2\subseteq[-L,L]^2\), and \(f_{P_\omega}\) is continuous on \(\mathcal X\) with
\[
L^{-1}\leq\frac1{48}\leq f_{P_\omega}(x)\leq\frac5{144}\leq L,
\qquad x\in\mathcal X .
\]
% lean: causalHardSquare_rectangularScoreSupport causalHardScoreDensity_continuous causalHardA1A2Law_density_certificates

The control regression is constant, and the treatment regression is the restriction to \(\mathcal X\) of the globally \(C^\infty\) function
\[
g_\omega(x)=\frac12+\frac{x^{(1)}}{16}
+\sum_{j=1}^{M}\omega_j\,\Delta\,\varphi\!\left(\frac{x-x_j}{w}\right),
\]
which satisfies the displayed envelope
\[
\max_{0\leq|\alpha|\leq p}\sup_{x\in\mathcal X}|\partial^\alpha g_\omega(x)|
+
\max_{0\leq|\alpha|\leq p}\sup_{x\neq z\in\mathcal X}
\frac{|\partial^\alpha g_\omega(x)-\partial^\alpha g_\omega(z)|}{\lVert x-z\rVert_2}
\leq 48\leq L .
\]
% lean: causalHardTreatmentProfile_has_smooth_extension causalHardControlProfile_euclideanCExtEnvelope causalHardTreatmentProfile_euclideanCExtEnvelope_powerBandwidth
The bump terms are controlled because the bandwidth is calibrated to the amplitude: for every \(0\leq k\leq q\) the \(k\)th derivative of one bump is bounded by \(\Delta w^{-k}b_k\), and
\[
\Delta\,w^{-k}b_k=\Delta^{1-k/q}A^{-k}b_k\leq1,
\]
since \(\Delta\leq1\) and \(k\leq q\) give \(\Delta^{1-k/q}\leq1\), while \(b_k\leq A\leq A^k\) for \(k\geq1\); at \(k=0\) the bound is \(\Delta b_0\leq\Delta A\leq1\).  The affine part \(1/2+x^{(1)}/16\) contributes at most \(1/2+3/16\) to the derivative maximum, because \(|x^{(1)}|/16\leq3/16\) on \(\mathcal X\) and its only nonvanishing derivative is the constant \(1/16\), and at most \(1/16\) to the Lipschitz maximum; and at most one bump is active at a time.
% lean: causalHardTreatmentProfile_powerBandwidth_derivative_scaling causalHardTreatmentProfile_euclideanCExtEnvelope_of_scaling_bounds

The conditional variances are the Bernoulli variances of the two means, so
\[
\sigma^2_{t,P_\omega}(x)=\mu_{t,P_\omega}(x)\{1-\mu_{t,P_\omega}(x)\}\in\bigl[\frac3{16},\frac14\bigr]
\subseteq[L^{-1},L],
\qquad x\in\mathcal X,\ t\in\{0,1\},
\]
and they are continuous on \(\mathcal X\).
% lean: causalSelectedBernoulliKernel_variance_id causalHardA1A2Law_kernel_certificates

Each \(Y(t)\) takes values in \(\{0,1\}\), so the moment envelope holds with room to spare:
\[
\mathbb E_{P_\omega}\bigl[|Y(t)|^{2+\nu}\mid X=x\bigr]\leq1\leq L,
\qquad x\in\mathcal X .
\]
% lean: causalSelectedBernoulliKernel_condAbsMoment_le_one causalHardA1A2Law_condAbsMoment_le

The sets \(\mathcal A_0,\mathcal A_1\) are a Borel partition of \(\mathcal X\) with \(\operatorname{bd}(\mathcal A_0)\cap\operatorname{bd}(\mathcal A_1)=\partial\mathcal A_1=\mathcal B\subset\operatorname{int}(\mathcal X)\), and \(\mathcal B\) is a compact rectifiable curve, being the boundary of a \(2\times2\) square.  Its size obeys
\[
2\leq\mathcal H^1(\mathcal B)\leq16,
\qquad\text{hence}\qquad
L^{-1}\leq\mathcal H^1(\mathcal B)\leq L,
\]
the lower bound because \(\mathcal B\) contains the segment from \((-1,0)\) to \((1,0)\) and the upper bound from a Lipschitz parametrization of the perimeter.
% lean: causalHardA1A2Law_class_geometry causalHardFrontier_hausdorff_bounds causalHardFrontier_hausdorff_bounds_fixed causalHardFrontier_rectifiableCurve

The metric is Euclidean and the kernel is \(K_\square\); the traces of closed Euclidean balls on \(\mathcal X\) form a VC class of index at most four by \cref{obj:lem:euclidean-balls-vc}.
% lean: euclidean_balls_vc

For the population Gram matrix, at every \(x\in\mathcal B\), both \(t\in\{0,1\}\), every \(0<h\leq L^{-1}\) and every \(v\in\mathbb R^{p+1}\),
\[
\frac1{48}\cdot\frac\pi2
\int_0^1\bigl\{v^\top r_p((2t-1)u)\bigr\}^2u\,du
\ \leq\
v^\top\Psi_{t,P_\omega,x}(h)\,v ,
\]
because the density is at least \(1/48\) and, in polar coordinates around \(x\), the arm-\(t\) part of the window \(B_2(x,h)\) contains a quarter-disk on which the uniform kernel equals one, namely a sector of opening \(\pi/2\) with one sign choice per coordinate determined by where \(x\) sits on \(\mathcal B\).  A quarter, not a half, is what is available at every interface point: at the four corners of \(\mathcal A_1\) the arm-\(t\) part of a small window is exactly a quarter-disk, and \(h\leq L^{-1}\leq1\) keeps that sector inside \(\mathcal X\).  Its angular measure \(\pi/2\) is the second factor in the display.  Combining this with the coercivity constant, \(\gamma_p\lVert v\rVert_2^2\leq v^\top\Psi_{t,P_\omega,x}(h)v\), and \(L\geq\gamma_p^{-1}\) gives the class threshold
\[
\lambda_{\min}\bigl(\Psi_{t,P_\omega,x}(h)\bigr)\geq L^{-1}.
\]
% lean: hardSquare_populationGram_quadratic_lower hardSquare_populationGramFloor_eventually causalHardA1A2Law_populationGram_certificate

For the small-bandwidth arm mass, at every \(x\in\mathcal B\), both \(t\), and every \(0<h\leq L^{-1}\),
\[
\frac1{h^2}\int_{\mathcal A_t}
K_\square\bigl(\{(2t-1)\lVert z-x\rVert_2\}/h\bigr)\,dz
=
\frac1{h^2}\int_{\mathcal A_t\cap B_2(x,h)}dz
\ \geq\
\frac18
\ \geq\
L^{-1},
\]
since every point of \(\partial\mathcal A_1\), corners included, retains at least an eighth of the disk area in each arm.
% lean: causalHardArm_inter_closedBall_volume_lower causalHardArmLocalMass_lower causalHardA1A2Law_localMass_certificate

Finally, for every \(x\in\mathcal B\), both \(t\), and every \(0<s\leq L^{-1}\), the one-dimensional Hausdorff integral of \(f_{P_\omega}\) over \(\{z\in\mathcal A_t:\lVert z-x\rVert_2=s\}\) is finite and strictly positive, because that slice is a nondegenerate circular arc and the density lies in \([1/48,5/144]\).
% lean: causalHardA1A2Law_slice_certificate
Hence \(P_\omega\in\mathcal P_{12}(p,\nu,L)\), with the common geometry displayed above.
% lean: causalHardA1A2Law_mem_class_of_analytic_certificates causalHardA1A2Law_geometry

\emph{Bit structure.}
Put \(\rho=\pi w^2/36>0\).  Point reflection through \(x_j\) maps the disk \(S_j\) onto itself, preserves \(\lVert x-x_j\rVert_2\) and negates \(\eta_j\), so the angular correction integrates to zero over \(S_j\); the remaining background density is the constant \(1/36\).  With the score projection of \cref{obj:synth_24},
\[
P_\omega\{z:\operatorname{score}(z)\in S_j\}
=
\frac1{36}\int_{S_j}dz
=
\frac{\pi w^2}{36}
=\rho .
\]
% lean: scoreCenter_closedBall_radial_angularTerm_integral_eq_zero causalHardScoreMeasure_cell_mass causalHardA1A2Law_cell_mass
Each bump and each angular correction indexed by \(k\) vanishes off \(S_k\), and the disks are disjoint; therefore the restriction of \(P_\omega\) to \(\{z:\operatorname{score}(z)\in S_j\}\) depends on \(\omega\) only through \(\omega_j\), and the restriction of \(P_\omega\) to the complement of \(\bigcup_{j}S_j\) is the same for every \(\omega\).
% lean: causalHardA1A2Law_restrict_cell_eq causalHardA1A2Law_restrict_off_cells_eq

\emph{Conditional signed-observation laws.}
Every \(z\in S_j\) has \(|z^{(1)}|\leq1/4+w\leq1\) and \(|z^{(2)}|\leq w\leq2\), so \(z\in\mathcal A_1\) exactly when \(z^{(2)}\geq0\).  Since \(x_j^{(2)}=0\), the signed distance of \cref{obj:def:cty-known-geometry-signed-distance-data} therefore reduces on \(S_j\) to the vertically signed radius,
\[
D^{\pm}_{P_\omega,x_j}
=
\begin{cases}
\ \ \lVert X-x_j\rVert_2, & X^{(2)}\geq 0,\\
-\lVert X-x_j\rVert_2, & X^{(2)}<0 .
\end{cases}
\]
% lean: causalHardSignedStatistic causalHardSignedStatistic_eq_verticalSignedRadius_on_cell
Using the conditional-law notation of \cref{obj:synth_25}, define
\[
Q_{j,b}
=
\mathcal L_{P_\omega}\bigl((Y,D^{\pm}_{P_\omega,x_j})\mid X\in S_j\bigr)
\qquad\text{for any }\omega\text{ with }\omega_j=b ,
\]
which is unambiguous by the one-bit locality just proved and is a probability law.  Equivalently, the image of \(P_\omega\) restricted to \(\{X\in S_j\}\) under \((Y,D^{\pm}_{P_\omega,x_j})\) equals \(\rho\,Q_{j,\omega_j}\).
% lean: causalHardCellBitObservationLaw causalHardCellBitRepresentative causalHardCellSignedObservationMeasure_eq_bitLaw

\emph{Signal separation.}
On \(\mathcal X\) the treatment-effect curve of \cref{obj:synth_46} is
\[
\tau_{P_\omega}(x)=\mu_{1,P_\omega}(x)-\mu_{0,P_\omega}(x)
=\frac{x^{(1)}}{16}+\sum_{j=1}^{M}\omega_j\Delta\,\varphi\!\left(\frac{x-x_j}{w}\right),
\]
and at \(x=x_j\) only the \(j\)th bump is active, with \(\varphi(0)=1\), so
\[
\tau_{P_\omega}(x_j)=\frac{x_j^{(1)}}{16}+\omega_j\Delta .
\]
% lean: packingRegression_center causalHardTreatmentProfile_center_flip
Hence two vertices agreeing in coordinate \(j\) have the same target at \(x_j\), and flipping coordinate \(j\) moves that target by exactly the amplitude,
\[
\bigl|\tau_{P_\omega}(x_j)-\tau_{P_{\omega^{(j)}}}(x_j)\bigr|=\Delta .
\]
% lean: causalHardA1A2Law_tau_center_eq_of_bit_eq causalHardA1A2Law_tau_center_flip

\emph{Common signed-distance marginal.}
The same reflection argument applies slice by slice: for every Borel \(A\subseteq\mathbb R\), the set \(\{z\in S_j:\ D^{\pm}_{P_\omega,x_j}(z)\in A\}\) is invariant under reflection through \(x_j\), so the angular correction contributes nothing to its mass and
\[
Q_{j,0}(\mathbb R\times A)
=
\frac1{36\rho}\int_{\{z\in S_j:\ D^{\pm}_{P_\omega,x_j}(z)\in A\}}dz
=
Q_{j,1}(\mathbb R\times A),
\]
which is the asserted equality of signed-distance marginals.
% lean: causalHardScoreMeasure_signedSlice_eq_uniform causalHardScoreMeasure_restrict_map_signedStatistic_eq causalHardCellBitObservationLaw_common_signedMarginal

\emph{Localized comparison of the two bit laws.}
Under \(Q_{j,b}\) the outcome is Bernoulli given the signed distance, with success parameter in \([1/4,3/4]\) by the middle-half bound above.  The two success parameters differ only at short positive radii: for every Borel \(B\subseteq\mathbb R\),
\[
\left|
\int_{\{u\in B\}} y\,dQ_{j,1}(y,u)
-
\int_{\{u\in B\}} y\,dQ_{j,0}(y,u)
\right|
\leq
\Delta\;Q_{j,0}\bigl(\mathbb R\times(B\cap(0,2C\Delta))\bigr).
\]
% lean: causalHardSignedSuccess_setIntegral_abs_le_cellArea causalHardSignedSuccess_localized_setwise_bound
Indeed, enabling the \(j\)th bit adds \(\Delta\varphi_{\mathrm r}(r/w)\) to the conditional mean on the upper half of the radial slice at radius \(r\) and simultaneously tilts the design density by \(t_\Delta(r)\eta_j\).  For \(r\geq2C\Delta=256\Delta\) the cutoff is fully active, \(\chi(r/(128\Delta)-1)=1\) and \(\max\{r/16,8\Delta\}=r/16\), so \(t_\Delta(r)=-32\Delta\varphi_{\mathrm r}(r/w)/r\); since the affine part of the treatment regression contributes the factor \((x^{(1)}-x_j^{(1)})/16=r\eta_j(x)/16\), the tilt contributes \(-2\Delta\varphi_{\mathrm r}(r/w)\eta_j(x)^2\) per unit slice mass, and the mean of \(\eta_j^2\) over the upper half circle of radius \(r\) is \(1/2\), which cancels the added bump exactly.  For \(0<r<256\Delta\) the two conditional means differ by at most \(\Delta\), and on the control side, where the signed distance is negative, both laws carry the common mean \(1/2\).
% lean: packingRegression_mul_density_enable_cell_radial_identity angularTilt angularCutoff_eq_one
Consequently the restrictions of the two laws to the complement of the short positive window agree: writing \(E=\{(y,u):0<u<2C\Delta\}\),
\[
Q_{j,0}(\,\cdot\,\cap E^{c})=Q_{j,1}(\,\cdot\,\cap E^{c}),
\]
which is the asserted agreement outside \(0<D^{\pm}_{P,x_j}<2C\Delta\).
% lean: statisticBernoulliOutcome_restrict_compl_eq_of_localized_success_bound causalHardCellSignedObservationMeasure_restrict_compl_enable_eq causalHardCellBitObservationLaw_restrict_compl_eq

\emph{The two divergence bounds.}
Two Bernoulli outcome laws over a common statistic, with success parameters in \([1/4,3/4]\) and setwise mean-mass discrepancy at most \(\Delta\) localized to \((0,2C\Delta)\), satisfy
\[
\operatorname{KL}(Q_{j,0},Q_{j,1})
\leq
4\Delta^2\,Q_{j,0}\bigl(\mathbb R\times(0,2C\Delta)\bigr),
\]
and the same bound with the two laws interchanged, the marginals being common.
% lean: statisticBernoulliOutcome_klDiv_le_of_localized_success_bound causalHardCellSignedObservationMeasure_klDiv_enable_le
Because \(\{z\in S_j:\ 0<D^{\pm}_{P_\omega,x_j}(z)<R\}\subseteq B_2(x_j,R)\), the localized mass obeys
\[
Q_{j,0}\bigl(\mathbb R\times(0,R)\bigr)
=
\frac1{36\rho}\int_{\{z\in S_j:\ 0<D^{\pm}_{P_\omega,x_j}(z)<R\}}dz
\leq
\frac{\pi R^2}{36\rho}.
\]
% lean: causalHardScoreMeasure_signedStatistic_Ioo_le
Taking \(R=2C\Delta=256\Delta\) and substituting \(\rho=\pi w^2/36\),
\[
4\Delta^2\cdot\frac{\pi(256\Delta)^2}{36\rho}
=
\frac{4\cdot 65536\,\Delta^4}{w^2}
=
C_0\frac{\Delta^4}{w^2},
\]
so both directed divergences are at most \(C_0\Delta^4/w^2\).
% lean: causalHardCellBitObservationLaw_klDiv_le

Collecting the displayed size and calibration identities, the common geometry, class membership, the bit structure, the conditional signed-observation laws, the exact signal separation, and the three local comparison facts gives every clause of the statement for the constants \(c,A,C=128,C_0=262144\) and the threshold \(\delta_0\).
% lean: causalHardHypercubeAt_of_constructorData A1A2HypercubeAt cty_a1_a2_rectangle_angular_hypercube
\end{proof}
